\chapter{Regression}\label{ch:iv1:regression}

\begin{ivtrains}
Regression as a quantitative researcher uses it: what a coefficient and its standard error really tell you, and what goes wrong on
financial data, where regressors are noisy, observations overlap, volatility moves and the signal is tiny. Most questions ask for
a number you can defend, and the later ones turn on an effect you could easily miss in your own work.
\end{ivtrains}

\begin{question}[$\star$ \iqroles{QR, QT, DS, Risk, Analyst} \iqfirm{any} \iqlevel{junior}]\label{iq:iv1:regression:1}
You regress $y$ on $x$ and get a slope of 2. What values can the slope of the reverse regression, $x$ on $y$, take? If the reverse slope turns out to be $0.18$, what is the correlation of $x$ and $y$?
\end{question}

\begin{question}[$\star$ \iqroles{QR, QT, DS, Risk} \iqfirm{any} \iqlevel{junior}]\label{iq:iv1:regression:2}
You copy every row of a dataset, making it twice as long, and refit the same linear regression. What happens to the coefficients, their standard errors, the $t$-statistics and $R^2$? Why does this matter when your observations overlap in time?
\end{question}

\begin{question}[$\star$ \iqroles{QR, QT, Trader, Risk} \iqfirm{any} \iqlevel{junior}]\label{iq:iv1:regression:3}
A stock has a volatility of 30\% a year and an index has 15\%, with correlation $0.8$. How many dollars of the index should you short per dollar of stock to make the hedged position's variance as small as possible? What volatility is left? What goes wrong if you take the ratio from the regression of the index on the stock, and invert it?
\end{question}

\begin{question}[$\star$ \iqroles{QR, DS, Risk, Analyst} \iqfirm{any} \iqlevel{junior}]\label{iq:iv1:regression:4}
A regression with an intercept reports $R^2$. What values can it take, and what does $R^2 = 1$ tell you? What happens to it when you add one more regressor? If you drop the intercept but still compute $1 - \text{RSS}/\sum (y_i - \bar y)^2$, where RSS is the sum of squared residuals, what can happen?
\end{question}

\begin{question}[$\star$ \iqroles{QR, QT, Analyst, Risk} \iqfirm{hedge fund, bank} \iqlevel{junior}]\label{iq:iv1:regression:5}
You regress a fund's returns on the market's and estimate an alpha, the regression's intercept, of 2\% a year. The residuals have a volatility of 10\% a year, and you have 4 years of data. What is the $t$-statistic of the alpha? If the true alpha is 2\%, how many years of data would you need for a $t$-statistic of 2?
\end{question}

\begin{question}[$\star$ \iqroles{QR, QT, DS} \iqfirm{hedge fund, market maker, crypto} \iqlevel{junior}]\label{iq:iv1:regression:6}
You regress daily returns on a standardised signal using 2\,500 days of data and get $R^2 = 0.5\%$. What is the $t$-statistic of the slope? Is the signal worth anything? Estimate the annual Sharpe ratio of holding a position proportional to the signal, where the Sharpe ratio is mean return divided by its standard deviation.
\end{question}

\begin{question}[$\star\star$ \iqroles{QR, QD, DS, Risk} \iqfirm{any} \iqlevel{junior}]\label{iq:iv1:regression:7}
Derive by hand the least squares estimate of $\beta$ in $y = X\beta + \varepsilon$. When does it exist? Why do you not compute it by inverting $X^\top X$? How do you fit it when $X$ is too large to fit in memory?
\end{question}

\begin{question}[$\star\star$ \iqroles{QR, QT, DS, Risk} \iqfirm{any} \iqlevel{junior, experienced}]\label{iq:iv1:regression:8}
You want the coefficient of $x$ in a regression of $y$ on $x$ and $z$. Procedure A regresses $y$ on $z$, and then regresses those residuals on $x$. Procedure B regresses $y$ on $z$ and $x$ on $z$, and then regresses the first set of residuals on the second. Which gives the same coefficient as the full regression? If $x$ and $z$ have correlation $0.8$ and the true coefficient of $x$ is 1, what does the other procedure give?
\end{question}

\begin{question}[$\star\star$ \iqroles{QR, QT, DS, Risk} \iqfirm{any} \iqlevel{junior}]\label{iq:iv1:regression:9}
The true model is $y = x - 3z + \varepsilon$, where $x$ and $z$ have unit variance and correlation $0.5$, and $\varepsilon$ is independent noise. You do not observe $z$, and you regress $y$ on $x$ alone. What slope do you get in a large sample?
\end{question}

\begin{question}[$\star\star$ \iqroles{QR, QT, DS, Risk} \iqfirm{any} \iqlevel{junior, experienced}]\label{iq:iv1:regression:10}
The true relation is $y = x^* + \varepsilon$. You observe $x = x^* + n$, where $x^*$ and the noise $n$ each have variance 1 and are independent. In a large sample, what slope do you get when you regress $y$ on $x$? What if the measurement noise is in $y$ and not in $x$? If both are noisy, how can you bracket the true slope with two regressions?
\end{question}

\begin{question}[$\star\star$ \iqroles{QR, DS, Risk} \iqfirm{any} \iqlevel{junior}]\label{iq:iv1:regression:11}
Two regressors have correlation $0.9$. By what factor is the standard error of each coefficient larger than if they were uncorrelated? What can you still estimate well? Can both $t$-statistics be below 2 while the $F$-test for the two together is highly significant?
\end{question}

\begin{question}[$\star\star$ \iqroles{QR, QT, DS} \iqfirm{any} \iqlevel{junior}]\label{iq:iv1:regression:12}
A point is chosen uniformly in the triangle with corners $(0,0)$, $(1,0)$ and $(1,1)$. Find the population regression line of $y$ on $x$, and the regression line of $x$ on $y$. What is the correlation of $x$ and $y$?
\end{question}

\begin{question}[$\star\star$ \iqroles{QR, DS, QD} \iqfirm{any} \iqlevel{junior, experienced}]\label{iq:iv1:regression:13}
You accidentally include the same column twice in a regression. What does ordinary least squares do? What does ridge do? What does lasso do?
\end{question}

\begin{question}[$\star\star$ \iqroles{QR, DS} \iqfirm{any} \iqlevel{junior, experienced}]\label{iq:iv1:regression:14}
The columns of $X$ are orthonormal, so that $X^\top X = I$, and the ordinary least squares coefficients are $3$, $-0.5$ and $1.2$. Find the coefficients that minimise $\tfrac12\lVert y - Xb\rVert^2 + \lambda\lVert b\rVert_1$ (lasso), and those that minimise $\tfrac12\lVert y - Xb\rVert^2 + \tfrac{\lambda}{2}\lVert b\rVert_2^2$ (ridge), with $\lambda = 1$. Derive the rule you use.
\end{question}

\begin{question}[$\star\star$ \iqroles{QR, QT, DS, Risk, Analyst} \iqfirm{any} \iqlevel{junior}]\label{iq:iv1:regression:15}
Under what assumptions are the coefficients of a linear regression unbiased, and under what further assumptions are their standard errors correct? On daily stock returns, which of these fail, and what do you do about each?
\end{question}

\begin{question}[$\star\star$ \iqroles{QR, QT, DS} \iqfirm{any} \iqlevel{junior}]\label{iq:iv1:regression:16}
You regress $y$ on 20 regressors with 100 observations and an intercept, and in truth $y$ is independent of all of them. What is the expected $R^2$? What is the expected adjusted $R^2$? How many of the 20 coefficients do you expect to have $|t| > 2$?
\end{question}

\begin{question}[$\star\star$ \iqroles{QR, QT, Risk, DS} \iqfirm{bank, hedge fund} \iqlevel{junior}]\label{iq:iv1:regression:17}
You generate two independent random walks, each 250 steps long, and regress one on the other, using the levels. How often is the slope ``significant'' with $|t| > 1.96$, and what $R^2$ do you see on average? What changes with 1\,000 steps? What should you do with price series instead?
\end{question}

\begin{question}[$\star\star$ \iqroles{QR, QT, DS} \iqfirm{hedge fund, market maker, bank} \iqlevel{junior, experienced}]\label{iq:iv1:regression:18}
You have 500 weak and correlated return signals and 5 years of daily data, about 1\,250 observations. Would you fit ridge or lasso? What would change your mind?
\end{question}

\begin{question}[$\star\star$ \iqroles{QR, QT, Risk, Trader} \iqfirm{bank, hedge fund} \iqlevel{junior, experienced}]\label{iq:iv1:regression:19}
A stock's true daily return is exactly $1.2$ times the market's. But it trades only on a random half of days, and on other days its price simply stays at the last trade. You regress its observed daily return on the same day's market return. What slope do you get? How do you recover $1.2$?
\end{question}

\begin{question}[$\star\star\star$ \iqroles{QR, DS, Risk} \iqfirm{any} \iqlevel{junior, experienced}]\label{iq:iv1:regression:20}
Let $x$ be standard normal and $y = \beta x + |x|\,\varepsilon$, with $\varepsilon$ standard normal and independent of $x$, so that the noise is large when $|x|$ is large. With a large sample of size $n$, what is the true standard error of the OLS slope, and what does the usual formula report? How do you fix the standard errors, and how could you use this to improve the estimate?
\end{question}

\begin{question}[$\star\star\star$ \iqroles{QR, QT, DS} \iqfirm{hedge fund, market maker} \iqlevel{junior, experienced}]\label{iq:iv1:regression:21}
Daily returns are independent, and a signal is unrelated to them. You regress the next 20 days' total return on today's signal, sampling every day. By what factor do the usual standard errors understate the truth if the signal is (a) independent from day to day, (b) nearly constant over 20 days, (c) an autoregressive process of order 1 with coefficient $0.9$?
\end{question}

\begin{question}[$\star\star\star$ \iqroles{QR, DS, QT} \iqfirm{any} \iqlevel{junior, experienced}]\label{iq:iv1:regression:22}
$y$ has correlation $0.6$ with $x_1$ and correlation $0.6$ with $x_2$, but you do not know the correlation between $x_1$ and $x_2$. What values can the $R^2$ of the regression of $y$ on both take? Then a different case: $x_1$ is uncorrelated with $y$, $x_2$ has correlation $0.4$ with $y$, and $x_1$ and $x_2$ have correlation $0.9$. Is $x_1$ useless?
\end{question}

\begin{question}[$\star\star\star$ \iqroles{QR, QT} \iqfirm{hedge fund, market maker, crypto} \iqlevel{experienced}]\label{iq:iv1:regression:23}
Two signals each have correlation $0.05$ with the next day's return. What correlation can the best linear combination of them reach if their correlation with each other is $0.5$, $0$, or $-0.5$? In the case of $0.5$, suppose the second signal's correlation with the return is $0.02$ and not $0.05$. What sign should its weight have?
\end{question}

\begin{question}[$\star\star\star$ \iqroles{QR, DS} \iqfirm{any} \iqlevel{experienced}]\label{iq:iv1:regression:24}
A regression has one regressor with $x^\top x = 1$, and noise of variance $\sigma^2 = 1$. The true coefficient is $\beta = 0.5$. OLS gives $\hat\beta$, and ridge gives $\hat\beta/(1 + \lambda)$. Which $\lambda$ minimises the mean squared error of the estimate, and how much lower is that error than for OLS? What does this say about shrinkage for weak signals?
\end{question}

\begin{question}[$\star\star\star$ \iqroles{QR, QT, Risk, DS} \iqfirm{any} \iqlevel{junior, experienced}]\label{iq:iv1:regression:25}
$x$ and $y$ are jointly normal with unit variance and correlation $0.5$. You keep only the observations with $x > 0$ and regress $y$ on $x$. What slope and correlation do you get? What slope do you get if instead you keep only the observations with $y > 0$?
\end{question}

\begin{question}[$\star\star\star$ \iqroles{QR, DS, Risk} \iqfirm{bank, hedge fund} \iqlevel{junior, experienced}]\label{iq:iv1:regression:26}
A regression of $y$ on $x$ with an intercept has 1\,000 observations, where $x$ has mean 0 and $\sum x_i^2 = 1000$. One observation has $x = 10$ and residual 2. Without refitting, find its leverage, the change in the slope if you remove it, and the prediction error for it from a fit without it.
\end{question}

\begin{question}[$\star\star\star$ \iqroles{QR, QT, DS} \iqfirm{hedge fund, bank, market maker} \iqlevel{experienced}]\label{iq:iv1:regression:27}
The true model is $y = x + u$, and $x = z + u + v$, where $z$, $u$ and $v$ are independent with unit variance. So $x$ is correlated with the error $u$. What slope does OLS of $y$ on $x$ give in a large sample? How do you use $z$ to get the true slope? What goes wrong if $z$ is only weakly related to $x$?
\end{question}

\begin{question}[$\star\star\star$ \iqroles{QR, QT, DS} \iqfirm{hedge fund, market maker} \iqlevel{junior, experienced}]\label{iq:iv1:regression:28}
You fit an OLS regression with 20 regressors on 1\,000 observations, and the true $R^2$ of the best linear predictor is 1\%. Roughly what $R^2$ should you expect out of sample, measured against predicting the true mean? How much data would you need for the out-of-sample $R^2$ to reach zero?
\end{question}

\begin{question}[$\star\star\star\star$ \iqroles{QR, QT} \iqfirm{hedge fund} \iqlevel{experienced}]\label{iq:iv1:regression:29}
Consider $r_{t+1} = \beta x_t + u_{t+1}$ and $x_{t+1} = \rho x_t + v_{t+1}$, where $\beta = 0$, $\rho = 0.95$, the shocks $u$ and $v$ have unit variance and correlation $-0.9$, and you have 120 observations. You regress $r_{t+1}$ on $x_t$ with an intercept. Roughly what is the mean of the estimated slope? You test $\beta = 0$ against $\beta > 0$ with the usual $t$-statistic at a nominal 5\% level. Does the test reject more or less often than 5\%, and roughly how often?
\end{question}

\begin{question}[$\star\star\star\star$ \iqroles{QR, QT, DS} \iqfirm{hedge fund, bank} \iqlevel{experienced}]\label{iq:iv1:regression:30}
$y = x_1 + \varepsilon$, so only $x_1$ matters. You observe $w = x_1 + n$, a noisy version of $x_1$ with noise variance $0.5$, along with a second variable $x_2$. Both $x_1$ and $x_2$ have unit variance, and their correlation is $0.8$, and $x_2$ has no effect on $y$ once $x_1$ is known. You regress $y$ on $w$ and $x_2$. What coefficients do you get in a large sample? What would you conclude about $x_2$, and what is the error?
\end{question}
